\section{Evaluation}
\label{sec:eval}

\subsection{Protocol}
\label{sec:eval:protocol}

We follow a frozen pipeline on branch \texttt{tube-v3-val}: code for \texttt{\_safe} and \texttt{\_fs}, pre-registered \(\alpha\) rule on validation (\texttt{configs/tube\_v3\_alpha\_rule.yaml}), freeze file \texttt{tube\_v3\_adaptive\_final\_envchunks.yaml}, and a single test run (\TubeEpisodes\ episodes per arm, master seed~\TubeMasterSeed, BBA host).
Primary arms are max-score \texttt{\_safe\_fs} with per-family \(\alpha\) from validation under \texttt{env\_chunks} calibration.
Confidence intervals use trace-cluster bootstrap (4000 resamples).

\subsection{Main results vs.\ RobustMPC}
\label{sec:eval:main}

\tabref{tab:tube:primary} summarizes the rule-selected primary arm.
On \mbox{5G} Markov and Irish traces, Tube-MPC substantially improves QoE and cuts rebuffering; joint coverage meets or exceeds \(1-\alpha\).
On harsh \mbox{5G}, rebuffering improves with QoE statistically similar to RobustMPC; fixed-\(\rho\) MPC still achieves higher QoE at roughly threefold rebuffering (full tables in the repository report).
On broadband, QoE gains are small and joint coverage (\TubeBBJoint) remains slightly below the nominal \(0.80\) for \(\alpha{=}0.20\), though chunk-aligned calibration raised joint coverage from \(0.646\) (raw-tick \texttt{\_fs}) to \TubeBBJoint; \(\alpha{=}0.15\) on the same test reaches \TubeBBJointA015.

\begin{cps@table}[!t]
\centering
\caption{Primary tube arm (\texttt{maxscore\_safe\_fs}, env\_chunks calibration) vs.\ RobustMPC on the frozen test split. Cluster 95\% CI.}
\label{tab:tube:primary}
\evaltable{%
\begin{tabular}{@{}lrrrc@{}}
\toprule
Trace set & QoE & Rebuf (s) & Joint cov. & $\Delta$QoE / $\Delta$rebuf (s) \\
\midrule
5g\_markov & \TubeMarkovQoE & \TubeMarkovRebuf & \TubeMarkovJoint &
  \cici{\TubeMarkovDQoE}{\TubeMarkovDQoELo}{\TubeMarkovDQoEHi} /
  \cici{\TubeMarkovDRebuf}{\TubeMarkovDRebufLo}{\TubeMarkovDRebufHi} \\
5g\_harsh & \TubeHarshQoE & \TubeHarshRebuf & \TubeHarshJoint &
  \cici{\TubeHarshDQoE}{\TubeHarshDQoELo}{\TubeHarshDQoEHi} /
  \cici{\TubeHarshDRebuf}{\TubeHarshDRebufLo}{\TubeHarshDRebufHi} \\
bb & \TubeBBQoE & \TubeBBRebuf & \TubeBBJoint &
  \cici{\TubeBBDQoE}{\TubeBBDQoELo}{\TubeBBDQoEHi} /
  \cici{\TubeBBDRebuf}{\TubeBBDRebufLo}{\TubeBBDRebufHi} \\
5g\_irish & \TubeIrishQoE & \TubeIrishRebuf & \TubeIrishJoint &
  \cici{\TubeIrishDQoE}{\TubeIrishDQoELo}{\TubeIrishDQoEHi} /
  \cici{\TubeIrishDRebuf}{\TubeIrishDRebufLo}{\TubeIrishDRebufHi} \\
\bottomrule
\end{tabular}%
}
\tablenote{Nominal joint targets: Markov/Irish \(1-\alpha\) with rule-selected \(\alpha\); harsh \(0.80\); bb \(0.80\) at \(\alpha{=}0.20\).}
\end{cps@table}

\begin{figure}[!t]
\centering
\evalfig{figures/pareto_qoe_rebuffer.png}
\caption{QoE vs.\ mean rebuffering on the frozen test split (placeholder path: copy from \texttt{docs/figures/tube\_v3\_adaptive\_envchunks/}).}
\label{fig:tube:pareto}
\end{figure}
